\chapter{Entropion and Ectropion Repair}

\section*{Introduction \& Clinical Indications}
Entropion and ectropion repair encompasses a series of surgical techniques aimed at correcting the malposition of the eyelids, specifically the inward turning (entropion) and outward turning (ectropion) of the eyelid margins. Entropion primarily affects the lower eyelid, leading to the eyelashes rubbing against the cornea and conjunctiva, which can result in irritation, tearing, and potential corneal damage. Conversely, ectropion results in the eyelid margin turning outward, exposing the conjunctiva and globe, causing dryness, irritation, and increased susceptibility to infections. These surgical interventions are critical for restoring the eyelid's normal anatomical position, alleviating symptoms, protecting the ocular surface, and enhancing both visual comfort and cosmetic appearance. The choice of surgical technique is influenced by the underlying causes, severity of the malposition, and specific anatomical considerations.

\begin{itemize}
  \item Eyelid tightening surgery
  \item Lower eyelid repair
  \item Tarsal strip procedure
  \item Lateral canthal tightening
  \item Everting suture repair
  \item Medial canthoplasty
  \item Full-thickness wedge resection
  \item Wies procedure
  \item Quickert-Rathbun suture repair
\end{itemize}

The primary surgical principle in entropion and ectropion repair is to restore the normal anatomical relationship of the eyelid margin to the globe, ensuring proper tear film distribution and protection of the ocular surface. In the case of entropion, the surgical goal is to evert the eyelid margin, preventing lash-cornea contact. This is typically achieved by addressing the underlying pathology, which often involves horizontal eyelid laxity, disinsertion or dehiscence of the lower eyelid retractors, and/or overriding of the preseptal orbicularis muscle. Techniques aim to tighten the eyelid horizontally, reattach or advance the lower eyelid retractors, and/or create scar tissue that pulls the eyelid margin outward, often by placing everting sutures through the orbicularis muscle and skin to rotate the eyelid margin outward. 

For ectropion, the primary objective is to tighten the horizontally lax eyelid and re-establish its apposition to the globe, allowing for effective tear drainage and protection from exposure. This is commonly accomplished by shortening the horizontal length of the eyelid, often through a lateral tarsal strip procedure where a portion of the lateral eyelid is resected and securely reattached to the lateral orbital rim periosteum. In cases of medial ectropion, a medial canthoplasty may be performed to tighten the medial canthal tendon and reposition the punctum. Cicatricial ectropion, caused by skin shortage or scarring, requires release of the scar tissue and provision of additional skin, often through skin grafting or local flap reconstruction. The choice of technique is meticulously tailored to address the specific anatomical defects contributing to the malposition, ensuring both functional and cosmetic restoration.

The surgical management of eyelid malposition has a rich history, with early references dating back to ancient Rome. Celsus, in the first century AD, described techniques involving skin excision to correct entropion and ectropion, reflecting an early understanding of the role of skin tension. Over the centuries, various methods were attempted, often with limited success, focusing on simple excisions, cautery, or lash epilation. 

The modern understanding of eyelid anatomy and the pathophysiology of involutional entropion and ectropion began to emerge in the late 19th and early 20th centuries, moving beyond empirical treatments to more anatomically informed approaches. Significant advancements in surgical techniques occurred in the mid-20th century. In the 1950s, Dr. Byron Smith introduced the "tarsal strip" concept for ectropion, which was later refined by Dr. Richard Tenzel in the 1970s, becoming the widely adopted lateral tarsal strip procedure. This technique revolutionized ectropion repair by effectively addressing horizontal eyelid laxity. 

For entropion, the Quickert-Rathbun everting suture technique, popularized in the 1970s by Dr. Marvin Quickert and Dr. William Rathbun, provided a simple yet effective method to correct lower eyelid retractor disinsertion and orbicularis override. Concurrently, the Wies procedure, involving a horizontal lid split and everting sutures, gained prominence for its efficacy. These innovations, combined with improved anesthesia, microsurgical instruments, and a deeper understanding of eyelid biomechanics, have transformed eyelid repair into a highly successful and predictable set of procedures, addressing not only functional deficits but also cosmetic concerns with durable results.

Surgical repair of entropion or ectropion is indicated when the eyelid malposition causes significant symptoms, threatens ocular health, or results in unacceptable cosmetic disfigurement.

\begin{itemize}
  \item Chronic foreign body sensation, irritation, or pain due to eyelashes rubbing against the cornea (trichiasis) in entropion.
  \item Corneal abrasion, ulceration, scarring, or infection resulting from chronic lash-cornea contact in entropion.
  \item Excessive tearing (epiphora) due to poor tear drainage or reflex tearing from ocular surface irritation in ectropion.
  \item Chronic conjunctivitis, recurrent ocular infections, or exposure keratopathy leading to corneal drying and damage in ectropion.
  \item Inability to fully close the eye (lagophthalmos) due to severe ectropion, increasing the risk of corneal desiccation.
  \item Significant cosmetic disfigurement affecting the patient's quality of life, self-esteem, or social interactions.
  \item Impaired vision secondary to corneal damage, chronic tearing obscuring the visual axis, or persistent ocular discomfort.
  \item Positive "snap-back" test or "distraction" test indicating significant horizontal eyelid laxity, particularly for ectropion.
  \item Evidence of lower eyelid retractor disinsertion or dehiscence on clinical examination, often contributing to involutional entropion.
\end{itemize}

Eyelid malposition repair is typically considered a primary surgical treatment for symptomatic and persistent entropion or ectropion. When conservative measures such as lubricating eye drops, temporary everting sutures (for entropion), or taping (for ectropion) fail to provide adequate relief or prevent ocular damage, surgery becomes the definitive first-line intervention. The decision to proceed with surgery is often driven by the risk of corneal damage, the severity of symptoms, and the impact on the patient's quality of life. 

For involutional forms, which constitute the majority of cases, surgery is the standard of care once symptoms become bothersome or vision is threatened. In cases where the malposition is recurrent after a previous surgical attempt, or if the condition is particularly complex due to extensive scarring (cicatricial), paralysis (e.g., facial nerve palsy), or congenital anomalies, the procedure may be considered secondary or revision surgery. Such secondary procedures often require more advanced reconstructive techniques, potentially involving skin grafts, local flaps, or specialized implants, to achieve a stable and functional outcome. The timing of surgery for paralytic ectropion may be delayed to allow for potential nerve recovery, making it a secondary intervention if conservative management fails during the waiting period.

\section*{Relevant Surgical Anatomy}
A comprehensive understanding of the eyelid anatomy is essential for successful entropion and ectropion repair. The eyelid is a multilayered structure consisting of skin, orbicularis oculi muscle, orbital septum, preaponeurotic fat pads, eyelid retractors (levator aponeurosis in the upper lid and capsulopalpebral fascia in the lower lid), tarsal plate, and conjunctiva. 

The orbicularis oculi muscle, innervated by the facial nerve (CN VII), plays a crucial role in eyelid closure and is divided into three parts: the pretarsal, preseptal, and orbital portions. The pretarsal orbicularis is particularly significant in the context of involutional entropion, where its overriding action can contribute to the condition. The lower eyelid retractors, originating from the inferior rectus muscle sheath, insert into the inferior border of the tarsal plate and the overlying skin, maintaining the eyelid's vertical position and apposition to the globe. 

The tarsal plates provide structural support and house the Meibomian glands, which are essential for tear film stability. The medial and lateral canthal tendons anchor the tarsal plates to the orbital rim, maintaining horizontal eyelid tension. Laxity of the lateral canthal tendon is a key factor in involutional ectropion. The lacrimal puncta, located on the medial aspect of the eyelid margins, are critical for tear drainage, and their malposition in ectropion leads to epiphora. 

Arterial supply to the eyelids is primarily from the medial and lateral palpebral arteries, branches of the ophthalmic artery, with contributions from the facial artery. Venous drainage parallels the arterial supply, emptying into the ophthalmic and facial veins. Sensory innervation is provided by branches of the trigeminal nerve (CN V). Lymphatic drainage from the lateral two-thirds of the eyelids drains to the preauricular lymph nodes, while the medial one-third drains to the submandibular nodes. Key surgical landmarks include the lateral orbital rim for canthal reattachment, the arcus marginalis, and the precise location of the lash line and gray line (mucocutaneous junction) for incisions and suture placement.

\section*{Preoperative Workup \& Preparation}
A thorough preoperative assessment is crucial for successful eyelid malposition repair. This typically includes a comprehensive ophthalmic examination, assessing visual acuity, ocular motility, and the integrity of the ocular surface, including fluorescein staining for corneal abrasions. A detailed evaluation of the eyelid anatomy is performed, including measurement of horizontal laxity (e.g., snap-back test, distraction test), assessment of lower eyelid retractor function, and identification of any underlying cicatricial components, such as skin shortage or vertical shortening. 

The patient's medical history is reviewed, particularly regarding any bleeding disorders, use of anticoagulants or antiplatelet medications, and systemic conditions that might affect healing or increase surgical risk.

\begin{itemize}
  \item \textbf{Medical Clearance:} Patients undergoing intravenous sedation or with significant medical comorbidities (e.g., cardiac disease, uncontrolled diabetes, severe pulmonary disease) may require medical clearance from their primary care physician or a cardiologist.

  \item \textbf{Laboratory Tests:} Routine preoperative blood tests, such as a complete blood count (CBC), basic metabolic panel (BMP), and coagulation studies (PT/INR, PTT), may be ordered, especially for patients on anticoagulants or with a history of bleeding.

  \item \textbf{Medication Review:} Patients are typically advised to discontinue blood-thinning medications (e.g., aspirin, NSAIDs, warfarin, clopidogrel) for a specified period (e.g., 5-10 days) before surgery, under the guidance of their prescribing physician, to minimize bleeding risk. Herbal supplements with anticoagulant properties (e.g., ginkgo biloba, ginseng, vitamin E) should also be stopped.

  \item \textbf{NPO Status:} Standard fasting guidelines (NPO for 6-8 hours for solids, 2 hours for clear liquids) are followed if intravenous sedation or general anesthesia is planned.

  \item \textbf{Informed Consent:} A detailed discussion of the procedure, potential risks (including undercorrection, overcorrection, recurrence, infection, bleeding, scarring, vision changes), expected benefits, and alternative treatments, followed by obtaining informed consent.

  \item \textbf{Prophylactic Antibiotics:} Topical antibiotic drops or ointment (e.g., erythromycin, tobramycin) may be prescribed for use before surgery to minimize infection risk, particularly if the ocular surface is compromised.

  \item \textbf{Photographs:} Preoperative photographs are routinely taken for documentation, surgical planning, and comparison with postoperative results.
\end{itemize}

\textbf{Absolute Contraindications:}
\begin{itemize}
  \item Active ocular or periocular infection: Any acute infection (e.g., bacterial conjunctivitis, blepharitis, orbital cellulitis, hordeolum) must be treated and completely resolved prior to surgery to prevent severe postoperative complications.

  \item Uncontrolled severe coagulopathy: Significant bleeding disorders that cannot be adequately managed preoperatively (e.g., severe hemophilia, critically low platelet count) pose an unacceptable risk of hemorrhage.

  \item Unrealistic patient expectations: Patients with expectations that cannot be realistically met by the surgical procedure, or those with body dysmorphic disorder, may be poor candidates.

  \item Patient refusal: The patient's informed decision to decline surgery.
\end{itemize}

\textbf{Relative Contraindications and Precautions:}
\begin{itemize}
  \item Severe dry eye syndrome: For ectropion repair, particularly the lateral tarsal strip, tightening the eyelid may exacerbate dry eye symptoms by reducing tear film distribution, requiring careful consideration, aggressive preoperative dry eye management, and patient counseling.

  \item Uncontrolled systemic diseases: Poorly controlled diabetes mellitus, hypertension, or other chronic medical conditions can increase surgical risks (e.g., infection, poor wound healing, cardiovascular events) and impair recovery. Optimization of these conditions is crucial.

  \item Previous eyelid surgery: Prior surgery can significantly alter normal anatomy, create scar tissue, and compromise tissue vascularity, making subsequent repairs more challenging and potentially increasing the risk of recurrence or complications.

  \item Patient non-compliance: Inability or unwillingness to follow detailed postoperative instructions (e.g., medication use, activity restrictions, wound care) can compromise healing and lead to suboptimal outcomes.

  \item Acute facial nerve palsy: In cases of paralytic ectropion, the timing of surgery may be delayed for 6-12 months to allow for potential spontaneous nerve recovery. Temporary measures like lubricating drops, taping, or gold weight implantation may be used in the interim.

  \item Thin skin or poor tissue quality: Patients with very thin, atrophic skin or compromised tissue quality may have a higher risk of wound dehiscence or less durable results.

  \item Glaucoma: While not a contraindication, patients with glaucoma require careful monitoring of intraocular pressure during and after surgery, especially if orbital manipulation is extensive.
\end{itemize}

\section*{Surgical Technique \& Procedure}
Eyelid malposition repair is most commonly performed under local anesthesia with intravenous sedation, ensuring patient comfort while allowing for intraoperative assessment of eyelid position and tension. The patient is positioned supine, and the surgical area is prepped with an antiseptic solution (e.g., povidone-iodine or chlorhexidine) and draped in a sterile fashion.

The specific steps vary significantly depending on whether entropion or ectropion is being corrected, and the underlying cause:

\begin{itemize}
  \item \textbf{Anesthesia and Incision:} Local anesthetic (e.g., 1\% or 2\% lidocaine with 1:100,000 epinephrine) is infiltrated into the eyelid to provide analgesia and vasoconstriction. For entropion, an incision may be made below the lash line (transcutaneous) or through the conjunctiva (transconjunctival). For ectropion, an incision is typically made at the lateral canthus or just below the lash line.

  \item \textbf{For Involutional Entropion (e.g., Quickert-Rathbun sutures):}
    \begin{itemize}
      \item Three double-armed 5-0 or 6-0 absorbable sutures (e.g., chromic gut or Vicryl) are placed full-thickness through the lower eyelid.
      \item Each suture starts from the inferior fornix, passes through the lower eyelid retractors and orbicularis muscle, and exits through the skin just below the lash line.
      \item When tied over a bolster or directly to the skin, these sutures evert the eyelid margin by tightening the retractors and creating a fibrotic scar that pulls the eyelid outward.
    \end{itemize}

  \item \textbf{For Involutional Entropion (e.g., Wies procedure):}
    \begin{itemize}
      \item A full-thickness horizontal incision is made through the tarsus and conjunctiva, approximately 4mm below the eyelid margin, extending for about two-thirds of the lid length.
      \item Three to four double-armed everting sutures are then placed, starting from the inferior edge of the tarsal incision, passing through the orbicularis and skin, and tied to rotate the distal tarsal plate outward, correcting the inward turn.
    \end{itemize}

  \item \textbf{For Involutional Ectropion (e.g., Lateral Tarsal Strip):}
    \begin{itemize}
      \item The lateral canthus is incised horizontally, and the inferior crus of the lateral canthal tendon is disinserted from the orbital rim.
      \item A strip of tarsus, approximately 3-5mm wide, is created by removing the mucocutaneous junction, conjunctiva, and orbicularis muscle from the lateral aspect of the lower eyelid.
      \item The tarsal strip is then shortened to the desired length, ensuring adequate tension, and securely reattached to the periosteum inside the lateral orbital rim, typically with a non-absorbable suture (e.g., 4-0 or 5-0 Mersilene or Prolene).
    \end{itemize}

  \item \textbf{For Medial Ectropion:}
    \begin{itemize}
      \item A diamond-shaped excision of conjunctiva and lower eyelid retractors may be performed medially, inferior to the punctum.
      \item Sutures are placed to tighten the medial canthal tendon and re-establish apposition of the punctum to the globe, often combined with a lateral tarsal strip.
    \end{itemize}

  \item \textbf{For Cicatricial Ectropion:}
    \begin{itemize}
      \item Scar tissue is released, often requiring a full-thickness skin incision below the lash line.
      \item A full-thickness skin graft (e.g., from the postauricular area or supraclavicular fossa) or a local skin flap is harvested and sutured into the defect to provide additional skin and counteract the cicatricial forces.
    \end{itemize}

  \item \textbf{Closure:} Incisions are meticulously closed with fine sutures (e.g., 6-0 or 7-0 Vicryl or nylon), which may be absorbable or non-absorbable depending on their location and the surgeon's preference.

  \item \textbf{Post-operative Care:} Antibiotic ointment (e.g., erythromycin or bacitracin) is applied to the surgical site, and a light dressing or eye shield may be placed to protect the area.
\end{itemize}

\section*{Complications \& Risks}
While entropion and ectropion repair are generally safe procedures with high success rates, potential risks and complications can occur. These can be broadly categorized as general surgical risks and specific procedure-related risks.

\textbf{General Surgical Risks:}
\begin{itemize}
  \item \textbf{Bleeding:} Bruising (ecchymosis) and swelling (edema) are common and expected. Significant hemorrhage or hematoma formation is rare but can occur, potentially requiring drainage.

  \item \textbf{Infection:} Although uncommon with proper sterile technique and prophylactic antibiotics, wound infection can occur, potentially requiring oral or intravenous antibiotic treatment or further intervention.

  \item \textbf{Anesthesia Complications:} Risks associated with local anesthesia and intravenous sedation include allergic reactions, respiratory depression, cardiovascular events (e.g., arrhythmia, hypotension), or nausea/vomiting, though these are rare.

  \item \textbf{Pain and Discomfort:} Postoperative pain is usually mild and manageable with over-the-counter analgesics, but persistent or severe pain should be evaluated.
\end{itemize}

\textbf{Procedure-Specific Risks:}
\begin{itemize}
  \item \textbf{Undercorrection or Recurrence:} The eyelid malposition may not be fully corrected, or it may recur over time, especially in cases with significant underlying laxity, cicatricial changes, or paralytic etiology. This may necessitate revision surgery.

  \item \textbf{Overcorrection:} For entropion repair, overcorrection can lead to iatrogenic ectropion (outward turning). For ectropion repair, overcorrection can result in lagophthalmos (inability to fully close the eye), leading to severe exposure keratopathy and corneal damage.

  \item \textbf{Asymmetry:} Minor differences in eyelid contour, height, or position between the two eyes may occur, which can be cosmetically noticeable.

  \item \textbf{Scarring:} While incisions are typically well-hidden along natural skin creases, visible scarring, hypertrophic scars, or keloids can occur, particularly in individuals predisposed to them.

  \item \textbf{Suture Granuloma or Extrusion:} Sutures, particularly non-absorbable ones, can sometimes cause irritation, form a granuloma (inflammatory nodule), or extrude through the skin or conjunctiva, requiring removal.

  \item \textbf{Corneal Damage:} Accidental injury to the cornea during surgery (e.g., with instruments or sutures) is rare but possible, potentially leading to abrasions or ulcers.

  \item \textbf{Persistent Tearing (Epiphora) or Dry Eye:} Postoperative tearing can persist if tear drainage is still impaired (e.g., punctal malposition) or if the ocular surface remains irritated. Conversely, dry eye symptoms may worsen if lid tightening is excessive, reducing tear film distribution.

  \item \textbf{Lash Abnormalities:} Misdirection, loss, or thinning of eyelashes (madarosis) can occur near the incision site.

  \item \textbf{Globe Perforation/Vision Loss:} Extremely rare but severe complications, such as globe perforation during suture placement or orbital hemorrhage leading to optic nerve compression, can potentially result in permanent vision loss.

  \item \textbf{Punctal Stenosis or Malposition:} Surgical manipulation near the medial canthus can lead to narrowing of the tear drainage puncta or their continued malposition, causing persistent epiphora.
\end{itemize}

\section*{Postoperative Care \& Recovery}
Immediately after the procedure, the patient will be monitored in the post-anesthesia care unit (PACU) until fully awake, stable, and pain is controlled. The surgical eye will typically have antibiotic ointment applied, and a light dressing or eye shield may be placed to protect the area from accidental rubbing or trauma. Patients are usually discharged home the same day with detailed written and verbal instructions for postoperative care.

During the first 24-48 hours, swelling, bruising (ecchymosis), and mild discomfort around the eye are common and expected. Cold compresses applied gently to the area for 15-20 minutes every hour while awake can significantly help reduce swelling and bruising. Patients are advised to keep their head elevated, even during sleep, and avoid strenuous activities, bending, or heavy lifting to minimize swelling and bleeding risk. Oral pain relievers (e.g., acetaminophen, ibuprofen) can manage any discomfort. Topical antibiotic drops or ointment are typically prescribed for several days to a week to prevent infection. Non-absorbable sutures, if used, are usually removed by the surgeon within 7 to 14 days post-surgery. 

Swelling and bruising gradually subside over the first 1-2 weeks, and most patients can resume light activities within a few days and return to normal activities, including work, within 1-2 weeks, avoiding activities that could put pressure on the eye or involve swimming. Long-term recovery involves the continued resolution of residual swelling, softening of scar tissue, and full maturation of the surgical site, which can take several months (typically 3-6 months) to achieve the final aesthetic and functional outcome.

During entropion and ectropion repair, the surgeon documents specific intraoperative findings that confirm the preoperative diagnosis and guide the surgical approach. For entropion, common findings include significant horizontal eyelid laxity (demonstrated by a positive distraction test), disinsertion or dehiscence of the lower eyelid retractors from the inferior tarsal border, and hypertrophy or overriding of the preseptal orbicularis muscle. For ectropion, findings typically include marked horizontal eyelid laxity (positive snap-back test), attenuation or disinsertion of the lateral canthal tendon, and in some cases, medial canthal tendon laxity or punctal eversion. Cicatricial cases will show evidence of skin shortage, vertical shortening, or dense scar tissue restricting eyelid movement.

Pathology reports are generally not applicable unless a biopsy was taken during the procedure for a suspicious lesion (e.g., a basal cell carcinoma causing cicatricial ectropion). In such instances, the pathology report would describe the histological characteristics of the excised tissue, confirming the diagnosis of the lesion and assessing surgical margins. For routine entropion or ectropion repair, no tissue is typically sent for pathological examination, as the "pathology" is a mechanical or structural defect rather than a cellular disease. Instead, the "results" are evaluated by the surgeon through follow-up clinical examinations, focusing on the functional and cosmetic improvements achieved.

A normal and successful postoperative course following entropion or ectropion repair is characterized by a progressive resolution of symptoms and restoration of normal eyelid function and appearance. Patients can expect initial swelling and bruising to gradually subside over 1-2 weeks. Mild pain or discomfort is typically well-controlled with over-the-counter analgesics and resolves within a few days. The primary goal is the stable correction of the eyelid malposition: for entropion, the eyelid margin should be everted with no lash-cornea contact; for ectropion, the eyelid should be properly apposed to the globe, with the punctum correctly positioned for tear drainage.

Patients should experience significant relief from preoperative symptoms such as foreign body sensation, irritation, excessive tearing, or dry eye. The ocular surface, if previously compromised, should show signs of healing, with resolution of corneal abrasions or ulcers. The final cosmetic appearance, while initially affected by swelling, should improve steadily over several weeks to months, resulting in a more natural and symmetrical eyelid contour. The ability to fully close the eye should be preserved or restored. Any persistent irritation, tearing, or signs of recurrence would indicate a less than optimal result, potentially requiring further intervention after adequate healing time.

Postoperative follow-up is essential to monitor healing, remove sutures, and assess the surgical outcome, ensuring long-term success and addressing any potential complications.

\begin{itemize}
  \item \textbf{Initial Postoperative Visit:} Typically scheduled for 1-3 days after surgery to check for early complications (e.g., hematoma, infection), assess swelling, evaluate eyelid position, and reinforce care instructions.

  \item \textbf{Suture Removal:} If non-absorbable skin sutures were used, they are usually removed at 7-14 days post-op, depending on the location and healing.

  \item \textbf{Subsequent Visits:} Further follow-up appointments may be scheduled at 1 month, 3 months, and possibly 6 months to monitor long-term healing, assess for recurrence, evaluate cosmetic outcome, and address any persistent concerns.

  \item \textbf{Eye Drops/Ointments:} Continued use of prescribed antibiotic or lubricating eye drops/ointments as directed by the surgeon, often for several weeks.

  \item \textbf{Activity Restrictions:} Patients are advised to avoid rubbing the eyes, strenuous activities, heavy lifting, swimming, and eye makeup for several weeks (typically 2-4 weeks) to prevent trauma and infection.

  \item \textbf{Warning Signs:} Patients are instructed to contact their surgeon immediately if they experience sudden severe pain, significant vision changes, excessive bleeding, pus-like discharge, worsening redness and swelling, or fever.

  \item \textbf{Scar Management:} Gentle massage of the incision lines may be recommended after sutures are removed and healing is underway to help soften scar tissue.

  \item \textbf{Revision Surgery:} In cases of undercorrection, overcorrection, or recurrence, revision surgery may be considered after adequate healing time, typically 3-6 months, to allow for tissue stabilization.
\end{itemize}

Adherence to these postoperative instructions is critical for optimal healing and a successful long-term outcome.

\section*{Outcomes \& Prognosis}
Entropion and ectropion are common eyelid malpositions, with their incidence increasing significantly with age. Involutional (age-related) forms are by far the most prevalent, accounting for the vast majority of cases in developed countries. Entropion is slightly more common than ectropion, with studies reporting prevalence rates of 2.1\% for entropion and 1.1\% for ectropion in individuals over 60 years of age in Western populations. Both conditions affect men and women, though some studies suggest a slight male predominance for entropion. 

While involutional forms are widespread globally, cicatricial entropion, particularly due to trachoma, remains a significant public health issue in endemic regions of Africa and Asia, affecting millions and leading to preventable blindness through chronic corneal damage. Congenital forms of both entropion and ectropion are rare. The morbidity associated with these conditions primarily stems from chronic ocular irritation, corneal damage (e.g., abrasions, ulcers, scarring), impaired vision, and persistent epiphora, all of which can significantly impact a patient's quality of life. Mortality directly attributable to entropion or ectropion is virtually non-existent, but the conditions can lead to severe visual impairment if left untreated.

The outcomes for surgical repair of involutional entropion and ectropion are generally excellent, with high success rates ranging from 80\% to 95\% for primary procedures. Most patients experience significant relief from symptoms, improved ocular comfort, and restoration of normal eyelid function and appearance. Functional outcomes include the preservation or improvement of corneal integrity, resolution of epiphora or dry eye symptoms, and improved visual comfort. 

The prognosis is generally very good, with long-lasting results, particularly for involutional forms where the underlying pathology (laxity, retractor disinsertion) is effectively addressed. However, recurrence rates vary depending on the specific technique used, the underlying etiology, and patient factors. Recurrence is more common in cicatricial or paralytic forms, or in cases with severe underlying tissue laxity that may progress over time. For instance, recurrence rates for involutional entropion are typically 5-10\%, while for cicatricial entropion, they can be higher, especially in trachoma-endemic areas. 

Prognostic factors influencing success include the surgeon's experience, accurate identification of the underlying cause of the malposition, and patient compliance with postoperative care. While a small percentage of patients may require revision surgery, the overall impact on quality of life and ocular health is overwhelmingly positive, making these highly effective and worthwhile surgical interventions. Although large-scale randomized controlled trials akin to those for systemic diseases are less common for these specific oculoplastic procedures, numerous prospective and retrospective studies and meta-analyses consistently demonstrate high efficacy and safety.

\section*{References}
\begin{enumerate}
  \item American Academy of Ophthalmology. Basic and Clinical Science Course, Section 7: Oculofacial Plastic and Orbital Surgery. American Academy of Ophthalmology; 2024.

  \item Kanski JJ, Bowling B. Kanski's Clinical Ophthalmology: A Systematic Approach. 9th ed. Elsevier; 2020.

  \item MedlinePlus. Ectropion and entropion repair. U.S. National Library of Medicine; 2024. Available from: \url{https://medlineplus.gov/ency/article/002936.htm}

  \item Yanoff M, Duker JS. Ophthalmology. 6th ed. Elsevier; 2023.

  \item Olver JM. Eyelid Surgery. In: Nduka C, editor. Principles and Practice of Ophthalmic Plastic and Reconstructive Surgery. Springer; 2017. p. 193-228.
\end{enumerate}

\clearpage
